%\include{pakage_tikz}
\usepackage{tikz}          % Charge le package TikZ
\usepackage{xcolor}        % Gestion avancée des couleurs
\usepackage{amsmath}       % Gestion des formules mathématiques
\usepackage{pgfplots}
\usepackage[active,tightpage]{preview}
%PREAMBULE pour schÃéma
\usepackage{pgfplots}
\usepackage{tikz}
\usepackage{listofitems}
\usepackage[european resistor, european voltage, european current]{circuitikz}
\usepackage{anyfontsize}
\usepackage{bm}
\usepackage[french, english]{babel}

\PreviewEnvironment{tikzpicture}
\pgfplotsset{compat=1.18}
\setlength\PreviewBorder{5pt}
\usetikzlibrary{arrows.meta,bending,positioning}
\usetikzlibrary{datavisualization.formats.functions}
\usetikzlibrary{spath3, intersections}
\usetikzlibrary{arrows,shapes,positioning}
\usetikzlibrary{decorations.text,decorations.markings,decorations.pathmorphing,
  decorations.pathreplacing}
\usetikzlibrary{calc,patterns,shapes.geometric}
\usetikzlibrary {fadings,patterns}
\usetikzlibrary{hobby} % <-- indispensable pour use Hobby shortcut
\usetikzlibrary{curvilinear}

%\usepgfplotslibrary{external}
%\tikzexternalize

%% Automatiser cela : une solution plus élégante  nouvelle page, tu peux utiliser la commande \newpage juste avant chaque \begin{tikzpicture}
%\let\oldtikzpicture\tikzpicture
%\let\endoldtikzpicture\endtikzpicture
%\renewenvironment{tikzpicture}{\newpage\oldtikzpicture}{\endoldtikzpicture}
%\usepackage{etoolbox} % dans le préambule, si pas encore utilisé
%\pretocmd{\tikzpicture}{\clearpage}{}{}

\AtBeginDocument{\shorthandoff{:;!?}}
\tikzset{every picture/.style={execute at begin picture={\shorthandoff{:;!?};}}}
\tikzstyle{every picture}+=[remember picture]
\tikzstyle{na} = [shape=rectangle,inner sep=0pt]
